%! Author = alejandrogonzalezturrubiates
%! Date = 18/08/25

\section*{Módulo 4: Matrices (placeholder)}
% Objetivos, operaciones básicas, propiedades.

\paragraph{Suma y resta de matrices}

\begin{ejercicio}
Calcule $A+B$, donde
\[
A=\begin{bmatrix} 2 & 3 \\ 4 & 5 \end{bmatrix}, \quad
B=\begin{bmatrix} 1 & 7 \\ 2 & 6 \end{bmatrix}.
\]
\end{ejercicio}
\muestraSolucion{
\[
A+B=\begin{bmatrix} 2+1 & 3+7 \\ 4+2 & 5+6 \end{bmatrix}=
\begin{bmatrix} 3 & 10 \\ 6 & 11 \end{bmatrix}.
\]
}

\begin{ejercicio}
Calcule $A-B$, donde
\[
A=\begin{bmatrix} 5 & 2 \\ 8 & 3 \end{bmatrix}, \quad
B=\begin{bmatrix} 1 & 4 \\ 6 & 2 \end{bmatrix}.
\]
\end{ejercicio}
\muestraSolucion{
\[
A-B=\begin{bmatrix} 5-1 & 2-4 \\ 8-6 & 3-2 \end{bmatrix}=
\begin{bmatrix} 4 & -2 \\ 2 & 1 \end{bmatrix}.
\]
}

\begin{ejercicio}
Calcule $A+B$, donde
\[
A=\begin{bmatrix} 3 & -5 & 2 \\ 7 & 0 & -4 \end{bmatrix}, \quad
B=\begin{bmatrix} -1 & 6 & 5 \\ 2 & 3 & 8 \end{bmatrix}.
\]
\end{ejercicio}
\muestraSolucion{
\[
A+B=\begin{bmatrix} 2 & 1 & 7 \\ 9 & 3 & 4 \end{bmatrix}.
\]
}

\begin{ejercicio}
Calcule $A-B$, donde
\[
A=\begin{bmatrix} 4 & 9 & -3 \\ 0 & 2 & 7 \end{bmatrix}, \quad
B=\begin{bmatrix} 1 & -3 & 5 \\ -4 & 6 & 2 \end{bmatrix}.
\]
\end{ejercicio}
\muestraSolucion{
\[
A-B=\begin{bmatrix} 3 & 12 & -8 \\ 4 & -4 & 5 \end{bmatrix}.
\]
}

\begin{ejercicio}
Calcule $A+B$, donde
\[
A=\begin{bmatrix} 1 & 2 & 3 \\ 4 & 5 & 6 \\ 7 & 8 & 9 \end{bmatrix}, \quad
B=\begin{bmatrix} 9 & 8 & 7 \\ 6 & 5 & 4 \\ 3 & 2 & 1 \end{bmatrix}.
\]
\end{ejercicio}
\muestraSolucion{
\[
A+B=\begin{bmatrix} 10 & 10 & 10 \\ 10 & 10 & 10 \\ 10 & 10 & 10 \end{bmatrix}.
\]
}

\begin{ejercicio}
Calcule $A-B$, donde
\[
A=\begin{bmatrix} 2 & 4 & 6 \\ 1 & 3 & 5 \\ 0 & -2 & -4 \end{bmatrix}, \quad
B=\begin{bmatrix} 1 & 1 & 1 \\ 1 & 1 & 1 \\ 1 & 1 & 1 \end{bmatrix}.
\]
\end{ejercicio}
\muestraSolucion{
\[
A-B=\begin{bmatrix} 1 & 3 & 5 \\ 0 & 2 & 4 \\ -1 & -3 & -5 \end{bmatrix}.
\]
}

\begin{ejercicio}
Calcule $A+B$, donde
\[
A=\begin{bmatrix} 3 & 2 \\ 5 & 7 \\ 1 & 4 \end{bmatrix}, \quad
B=\begin{bmatrix} 4 & 1 \\ 2 & 3 \\ 6 & 5 \end{bmatrix}.
\]
\end{ejercicio}
\muestraSolucion{
\[
A+B=\begin{bmatrix} 7 & 3 \\ 7 & 10 \\ 7 & 9 \end{bmatrix}.
\]
}

\begin{ejercicio}
Calcule $A-B$, donde
\[
A=\begin{bmatrix} 8 & 5 \\ 3 & 9 \\ 6 & 7 \end{bmatrix}, \quad
B=\begin{bmatrix} 4 & 2 \\ 1 & 8 \\ 2 & 5 \end{bmatrix}.
\]
\end{ejercicio}
\muestraSolucion{
\[
A-B=\begin{bmatrix} 4 & 3 \\ 2 & 1 \\ 4 & 2 \end{bmatrix}.
\]
}

\begin{ejercicio}
Calcule $A+B$, donde
\[
A=\begin{bmatrix} 1 & -1 \\ -2 & 2 \end{bmatrix}, \quad
B=\begin{bmatrix} 3 & 3 \\ 4 & 4 \end{bmatrix}.
\]
\end{ejercicio}
\muestraSolucion{
\[
A+B=\begin{bmatrix} 4 & 2 \\ 2 & 6 \end{bmatrix}.
\]
}

\begin{ejercicio}
Calcule $A-B$, donde
\[
A=\begin{bmatrix} 5 & 8 & -3 \\ 0 & -1 & 2 \end{bmatrix}, \quad
B=\begin{bmatrix} 2 & -4 & 1 \\ -3 & 5 & -2 \end{bmatrix}.
\]
\end{ejercicio}
\muestraSolucion{
\[
A-B=\begin{bmatrix} 3 & 12 & -4 \\ 3 & -6 & 4 \end{bmatrix}.
\]
}

%! Author = alejandrogonzalezturrubiates
%! Date = 01/10/25

% ===== Ejercicio 11 =====
\begin{ejercicio}
Calcule $A+B$, donde
\[
A=\begin{bmatrix} 4 & 6 \\ 2 & 8 \end{bmatrix}, \quad
B=\begin{bmatrix} 1 & 3 \\ 5 & 7 \end{bmatrix}.
\]
\end{ejercicio}
\muestraSolucion{
\[
A+B=\begin{bmatrix} 5 & 9 \\ 7 & 15 \end{bmatrix}.
\]
}

% ===== Ejercicio 12 =====
\begin{ejercicio}
Calcule $A-B$, donde
\[
A=\begin{bmatrix} 10 & 4 \\ 6 & 9 \end{bmatrix}, \quad
B=\begin{bmatrix} 3 & 2 \\ 4 & 1 \end{bmatrix}.
\]
\end{ejercicio}
\muestraSolucion{
\[
A-B=\begin{bmatrix} 7 & 2 \\ 2 & 8 \end{bmatrix}.
\]
}

% ===== Ejercicio 13 =====
\begin{ejercicio}
Calcule $A+B$, donde
\[
A=\begin{bmatrix} 1 & 2 & 3 \\ 4 & 5 & 6 \end{bmatrix}, \quad
B=\begin{bmatrix} 6 & 5 & 4 \\ 3 & 2 & 1 \end{bmatrix}.
\]
\end{ejercicio}
\muestraSolucion{
\[
A+B=\begin{bmatrix} 7 & 7 & 7 \\ 7 & 7 & 7 \end{bmatrix}.
\]
}

% ===== Ejercicio 14 =====
\begin{ejercicio}
Calcule $A-B$, donde
\[
A=\begin{bmatrix} 7 & 9 \\ 2 & 5 \end{bmatrix}, \quad
B=\begin{bmatrix} 4 & 3 \\ 1 & 7 \end{bmatrix}.
\]
\end{ejercicio}
\muestraSolucion{
\[
A-B=\begin{bmatrix} 3 & 6 \\ 1 & -2 \end{bmatrix}.
\]
}

% ===== Ejercicio 15 =====
\begin{ejercicio}
Calcule $A+B$, donde
\[
A=\begin{bmatrix} 0 & 1 & -2 \\ 3 & 4 & 5 \end{bmatrix}, \quad
B=\begin{bmatrix} 6 & -1 & 3 \\ -3 & 2 & 7 \end{bmatrix}.
\]
\end{ejercicio}
\muestraSolucion{
\[
A+B=\begin{bmatrix} 6 & 0 & 1 \\ 0 & 6 & 12 \end{bmatrix}.
\]
}

% ===== Ejercicio 16 =====
\begin{ejercicio}
Calcule $A-B$, donde
\[
A=\begin{bmatrix} 2 & 4 & 6 \\ 8 & 10 & 12 \end{bmatrix}, \quad
B=\begin{bmatrix} 1 & 3 & 5 \\ 7 & 9 & 11 \end{bmatrix}.
\]
\end{ejercicio}
\muestraSolucion{
\[
A-B=\begin{bmatrix} 1 & 1 & 1 \\ 1 & 1 & 1 \end{bmatrix}.
\]
}

% ===== Ejercicio 17 =====
\begin{ejercicio}
Calcule $A+B$, donde
\[
A=\begin{bmatrix} 3 & -5 \\ 2 & 7 \\ 6 & 1 \end{bmatrix}, \quad
B=\begin{bmatrix} -1 & 8 \\ 4 & -2 \\ 3 & 5 \end{bmatrix}.
\]
\end{ejercicio}
\muestraSolucion{
\[
A+B=\begin{bmatrix} 2 & 3 \\ 6 & 5 \\ 9 & 6 \end{bmatrix}.
\]
}

% ===== Ejercicio 18 =====
\begin{ejercicio}
Calcule $A-B$, donde
\[
A=\begin{bmatrix} 9 & 4 \\ 7 & 8 \\ 5 & 6 \end{bmatrix}, \quad
B=\begin{bmatrix} 2 & 1 \\ 3 & 5 \\ 4 & 2 \end{bmatrix}.
\]
\end{ejercicio}
\muestraSolucion{
\[
A-B=\begin{bmatrix} 7 & 3 \\ 4 & 3 \\ 1 & 4 \end{bmatrix}.
\]
}

% ===== Ejercicio 19 =====
\begin{ejercicio}
Calcule $A+B$, donde
\[
A=\begin{bmatrix} -3 & 2 \\ 4 & -5 \end{bmatrix}, \quad
B=\begin{bmatrix} 6 & 1 \\ -2 & 8 \end{bmatrix}.
\]
\end{ejercicio}
\muestraSolucion{
\[
A+B=\begin{bmatrix} 3 & 3 \\ 2 & 3 \end{bmatrix}.
\]
}

% ===== Ejercicio 20 =====
\begin{ejercicio}
Calcule $A-B$, donde
\[
A=\begin{bmatrix} 8 & 5 & 2 \\ 7 & 6 & 1 \\ 4 & 3 & 0 \end{bmatrix}, \quad
B=\begin{bmatrix} 4 & 2 & 1 \\ 1 & 3 & 5 \\ 2 & 4 & 6 \end{bmatrix}.
\]
\end{ejercicio}
\muestraSolucion{
\[
A-B=\begin{bmatrix} 4 & 3 & 1 \\ 6 & 3 & -4 \\ 2 & -1 & -6 \end{bmatrix}.
\]
}

% ===== Ejercicio 21 =====
\begin{ejercicio}
Calcule $A+B$, donde
\[
A=\begin{bmatrix} 2 & 3 & 4 \\ 0 & -1 & 5 \\ 6 & 7 & 8 \end{bmatrix}, \quad
B=\begin{bmatrix} 1 & 0 & -2 \\ 3 & 4 & -1 \\ 5 & 6 & 7 \end{bmatrix}.
\]
\end{ejercicio}
\muestraSolucion{
\[
A+B=\begin{bmatrix} 3 & 3 & 2 \\ 3 & 3 & 4 \\ 11 & 13 & 15 \end{bmatrix}.
\]
}

% ===== Ejercicio 22 =====
\begin{ejercicio}
Calcule $A-B$, donde
\[
A=\begin{bmatrix} 1 & 2 \\ 3 & 4 \end{bmatrix}, \quad
B=\begin{bmatrix} 4 & 3 \\ 2 & 1 \end{bmatrix}.
\]
\end{ejercicio}
\muestraSolucion{
\[
A-B=\begin{bmatrix} -3 & -1 \\ 1 & 3 \end{bmatrix}.
\]
}

% ===== Ejercicio 23 =====
\begin{ejercicio}
Calcule $A+B$, donde
\[
A=\begin{bmatrix} 5 & -3 & 2 \\ 1 & 4 & -1 \end{bmatrix}, \quad
B=\begin{bmatrix} -2 & 6 & 1 \\ 3 & -1 & 2 \end{bmatrix}.
\]
\end{ejercicio}
\muestraSolucion{
\[
A+B=\begin{bmatrix} 3 & 3 & 3 \\ 4 & 3 & 1 \end{bmatrix}.
\]
}

% ===== Ejercicio 24 =====
\begin{ejercicio}
Calcule $A-B$, donde
\[
A=\begin{bmatrix} 7 & 8 \\ 6 & 5 \end{bmatrix}, \quad
B=\begin{bmatrix} 2 & 4 \\ 3 & 7 \end{bmatrix}.
\]
\end{ejercicio}
\muestraSolucion{
\[
A-B=\begin{bmatrix} 5 & 4 \\ 3 & -2 \end{bmatrix}.
\]
}

% ===== Ejercicio 25 =====
\begin{ejercicio}
Calcule $A+B$, donde
\[
A=\begin{bmatrix} 9 & 1 & 3 \\ 4 & 0 & 2 \\ 8 & 5 & 7 \end{bmatrix}, \quad
B=\begin{bmatrix} 1 & 3 & 2 \\ 5 & 7 & 8 \\ 4 & 6 & 9 \end{bmatrix}.
\]
\end{ejercicio}
\muestraSolucion{
\[
A+B=\begin{bmatrix} 10 & 4 & 5 \\ 9 & 7 & 10 \\ 12 & 11 & 16 \end{bmatrix}.
\]
}

% ===== Ejercicio 26 =====
\begin{ejercicio}
Calcule $A-B$, donde
\[
A=\begin{bmatrix} 3 & 4 & 5 \\ 2 & 3 & 4 \\ 1 & 2 & 3 \end{bmatrix}, \quad
B=\begin{bmatrix} 1 & 2 & 3 \\ 3 & 4 & 5 \\ 5 & 6 & 7 \end{bmatrix}.
\]
\end{ejercicio}
\muestraSolucion{
\[
A-B=\begin{bmatrix} 2 & 2 & 2 \\ -1 & -1 & -1 \\ -4 & -4 & -4 \end{bmatrix}.
\]
}

% ===== Ejercicio 27 =====
\begin{ejercicio}
Calcule $A+B$, donde
\[
A=\begin{bmatrix} 2 & 5 \\ -3 & 4 \end{bmatrix}, \quad
B=\begin{bmatrix} 1 & -2 \\ 3 & -4 \end{bmatrix}.
\]
\end{ejercicio}
\muestraSolucion{
\[
A+B=\begin{bmatrix} 3 & 3 \\ 0 & 0 \end{bmatrix}.
\]
}

% ===== Ejercicio 28 =====
\begin{ejercicio}
Calcule $A-B$, donde
\[
A=\begin{bmatrix} -4 & 2 & 7 \\ 5 & -3 & 0 \end{bmatrix}, \quad
B=\begin{bmatrix} 2 & 5 & -3 \\ -4 & 2 & 6 \end{bmatrix}.
\]
\end{ejercicio}
\muestraSolucion{
\[
A-B=\begin{bmatrix} -6 & -3 & 10 \\ 9 & -5 & -6 \end{bmatrix}.
\]
}

% ===== Ejercicio 29 =====
\begin{ejercicio}
Calcule $A+B$, donde
\[
A=\begin{bmatrix} 1 & 1 & 1 \\ 2 & 2 & 2 \\ 3 & 3 & 3 \end{bmatrix}, \quad
B=\begin{bmatrix} 3 & 2 & 1 \\ 1 & 2 & 3 \\ 3 & 1 & 2 \end{bmatrix}.
\]
\end{ejercicio}
\muestraSolucion{
\[
A+B=\begin{bmatrix} 4 & 3 & 2 \\ 3 & 4 & 5 \\ 6 & 4 & 5 \end{bmatrix}.
\]
}

% ===== Ejercicio 30 =====
\begin{ejercicio}
Calcule $A-B$, donde
\[
A=\begin{bmatrix} 10 & 9 & 8 \\ 7 & 6 & 5 \\ 4 & 3 & 2 \end{bmatrix}, \quad
B=\begin{bmatrix} 1 & 2 & 3 \\ 4 & 5 & 6 \\ 7 & 8 & 9 \end{bmatrix}.
\]
\end{ejercicio}
\muestraSolucion{
\[
A-B=\begin{bmatrix} 9 & 7 & 5 \\ 3 & 1 & -1 \\ -3 & -5 & -7 \end{bmatrix}.
\]
}